\medskip

\textbf{Partie A : dénombrement}

On considère l'ensemble des nombres entiers relatifs \textbf{non nuls} compris entre $- 30$ et $30$ ; cet ensemble peut s'écrire ainsi : \{$- 30~;~- 29~;~- 28~;~ ... - 1~;~1~;~ ... ~;~ 28~;~29~;~30$\}. Il comporte 60 éléments.

On choisit dans cet ensemble successivement et sans remise un entier relatif $a$ puis un entier relatif $c$.

\medskip

\begin{enumerate}
\item Combien de couples $(a~;~ c)$ différents peut-on ainsi obtenir ?
\end{enumerate}

On considère l'évènement $M$ : \og l'équation $ax^2 + 2x + c = 0$ possède deux solutions réelles distinctes \fg, où $a$ et $c$ sont les entiers relatifs précédemment choisis.

\begin{enumerate}[resume]
\item Montrer que l'évènement $M$ a lieu si et seulement si $ac < 1$.
\item Expliquer pourquoi l'évènement contraire $\overline{M}$ comporte \np{1740} issues.
\item Quelle est la probabilité de l'évènement $M$ ? On arrondira le résultat à $10^{-2}$.
\end{enumerate}

\medskip

\textbf{Partie B : équation différentielle}

\medskip

On considère l'équation différentielle 

\begin{center}$(E) :\quad  y' + 10y = \left(30x^2 + 22x - 8\right)\e^{-5x+1}$\quad avec \quad $x \in \R$\end{center}

où $y$ est une fonction définie et dérivable sur $\R$.

\medskip

\begin{enumerate}
\item Résoudre sur $\R$ l'équation différentielle : $y' + 10y = 0$.
\item Soit la fonction $f$ définie sur $\R$ par 

\[f(x) = (6x^2 + 2x - 2)\e^{-5x+ 1}.\]

On admet que $f$ est dérivable sur $\R$ et on note $f'$ la fonction dérivée de la fonction $f$.

Justifier que $f$ est une solution particulière de $(E)$.
\item Donner l'expression de toutes les solutions de $(E)$.
\end{enumerate}

\medskip

\textbf{Partie C : étude de fonction}

\medskip

On propose d'étudier dans cette partie la fonction $f$ rencontrée à la partie B question 2. 

On rappelle que, pour tout réel $x\:, f(x) = \left(6x^2 + 2x - 2\right)\e^{-5x+ 1}$.

On note $f'$ la fonction dérivée de la fonction $f$. On appelle $\mathcal{C}_f$ la courbe représentative de $f$ dans un repère du plan.

\medskip

\begin{enumerate}
\item On admet que $\displaystyle\lim_{x \to + \infty} f(x) = 0$.

Déterminer la limite de la fonction $f$ en $- \infty$
\item En utilisant la partie A, montrer que $\mathcal{C}_f$ coupe l'axe des abscisses en deux points (les coordonnées de ces points ne sont pas attendues).
\item En utilisant les parties A et B, montrer que $\mathcal{C}_f$ possède deux tangentes horizontales.
\item Dresser le tableau de variation complet de la fonction $f$.
\item Déterminer en justifiant le nombre de solution(s) de l'équation $f(x)= 1$.
\item Pour tout réel $m$ strictement supérieur à $0,2$, on définit $I_m$ par $I_m = \displaystyle\int_{0,2}^m f(x)\:\text{d}x$.
	\begin{enumerate}
		\item Vérifier que la fonction $F$ définie sur $\R$ par

\[F(x) = \left(-\dfrac65 x^2 - \dfrac{22}{25}x + \dfrac{28}{125}\right)\e^{-5x+ 1}\]
est une primitive de la fonction $f$ sur $\R$.
		\item Existe-t-il une valeur de $m$ pour laquelle $I_m = 0$ ? 

Interpréter graphiquement ce résultat.
	\end{enumerate}
\end{enumerate}
